% GENERATED FILE. Edit canonical lessons, then run npm run generate:books.
% canonical-source-hash: fnv1a64:55dabac38023172d
% canonical-lessons: GU-C438-ujas, GU-C438-sandhya, GU-C438-jhakhu, GU-C438-dhudhlu, GU-C438-tamtamvu

\chapter{Dawn and Dusk}
\label{ch:gu-dawn-and-dusk}
\hlchaptermodality{438}

\begin{chapteropening}
\textbf{By the end of this chapter:} \emph{I can talk about the light at dawn and at dusk.}

Complete the last lesson of chapter 438: I can talk about the light at dawn and at dusk.
\end{chapteropening}

\section[\texorpdfstring{\gu{ઉજાસ} (ujās)}{ઉજાસ (ujās)}]{\gu{ઉજાસ} (ujās) — the light of daybreak; brightness}

\label{lesson:GU-C438-ujas}



\begin{quote}
Before the new one: say the Gujarati for to congeal, then the Gujarati for to shiver all over.
\end{quote}

\subsection*{You'll want to know: \gu{ઉજાસ}}
\textbf{\gu{ઉજાસ}} — \emph{ujās} — \textquotedblleft{}the light of daybreak; brightness\textquotedblright{}.

\begin{quote}
\gu{સવારે} \gu{બારીમાંથી} \gu{ઓરડામાં} \gu{ઉજાસ} \gu{આવ્યો}.
\end{quote}

\begin{quote}
\emph{savāre bārīmā̃thī orḍāmā̃ ujās āvyo.}
\end{quote}

\begin{quote}
In the morning light came into the room through the window.
\end{quote}

A masculine noun.

\begin{cousinweb}
\textbf{\gu{ઉજાસ}} (\emph{ujās}) goes back to \emph{*ujjhāṣa}, \textquotedblleft{}shining\textquotedblright{}, a form CDIAL reconstructs; CDIAL lists the Gujarati word, \textquotedblleft{}brightness, light\textquotedblright{} (CDIAL 1678).
\end{cousinweb}

\begin{grammarlens}[title={What you've built}]
One more word for getting around.
\end{grammarlens}

\subsection*{Your turn}
\begin{itemize}
\raggedright
  \item \emph{Say it:} \emph{ujās}
  \item \emph{Say it:} \emph{ujās}, once more
  \item \emph{Say it:} \emph{savāre bārīmā̃thī orḍāmā̃ ujās āvyo.} — In the morning light came into the room through the window.
  \item \emph{Recall:} say the Gujarati for to congeal, then the Gujarati for to shiver all over, then say \emph{ujās} again
\end{itemize}

\begin{tcolorbox}[breakable,colback=teal!4,colframe=teal!35!black,title={Before you move on}]
What does \gu{ઉજાસ} mean? (\textquotedblleft{}The light of daybreak; brightness\textquotedblright{}.) Say it once more.
\end{tcolorbox}
\section[\texorpdfstring{\gu{સંધ્યા} (sandhyā)}{સંધ્યા (sandhyā)}]{\gu{સંધ્યા} (sandhyā) — dusk; the evening twilight}

\label{lesson:GU-C438-sandhya}



\begin{quote}
Before the new one: say the Gujarati for to shiver all over, then the Gujarati for the light of daybreak.
\end{quote}

\subsection*{You'll want to know: \gu{સંધ્યા}}
\textbf{\gu{સંધ્યા}} — \emph{sandhyā} — \textquotedblleft{}dusk; the evening twilight\textquotedblright{}.

\begin{quote}
\gu{સંધ્યાનું} \gu{આકાશ} \gu{ગુલાબી} \gu{છે}.
\end{quote}

\begin{quote}
\emph{sandhyānũ ākāś gulābī che.}
\end{quote}

\begin{quote}
The sky at dusk is pink.
\end{quote}

A feminine noun.

\begin{cousinweb}
\textbf{\gu{સંધ્યા}} (\emph{sandhyā}) is a tatsama, Sanskrit \emph{saṁdhyā}, \textquotedblleft{}a joining\textquotedblright{}, the meeting of day and night (Monier-Williams).
\end{cousinweb}

\begin{grammarlens}[title={What you've built}]
One more word for getting around.
\end{grammarlens}

\subsection*{Your turn}
\begin{itemize}
\raggedright
  \item \emph{Say it:} \emph{sandhyā}
  \item \emph{Say it:} \emph{sandhyā}, once more
  \item \emph{Say it:} \emph{sandhyānũ ākāś gulābī che.} — The sky at dusk is pink.
  \item \emph{Recall:} say the Gujarati for to shiver all over, then the Gujarati for the light of daybreak, then say \emph{sandhyā} again
\end{itemize}

\begin{tcolorbox}[breakable,colback=teal!4,colframe=teal!35!black,title={Before you move on}]
What does \gu{સંધ્યા} mean? (\textquotedblleft{}Dusk; the evening twilight\textquotedblright{}.) Say it once more.
\end{tcolorbox}
\section[\texorpdfstring{\gu{ઝાંખું} (jhā̃khũ)}{ઝાંખું (jhā̃khũ)}]{\gu{ઝાંખું} (jhā̃khũ) — dim; faint (of light)}

\label{lesson:GU-C438-jhakhu}



\begin{quote}
Before the new one: say the Gujarati for the light of daybreak, then the Gujarati for dusk.
\end{quote}

\subsection*{You'll want to know: \gu{ઝાંખું}}
\textbf{\gu{ઝાંખું}} — \emph{jhā̃khũ} — \textquotedblleft{}dim; faint (of light)\textquotedblright{}.

\begin{quote}
\gu{ઓરડામાં} \gu{ઝાંખી} \gu{બત્તી} \gu{હતી}.
\end{quote}

\begin{quote}
\emph{orḍāmā̃ jhā̃khī battī hatī.}
\end{quote}

\begin{quote}
There was a dim light in the room.
\end{quote}

\begin{cousinweb}
\textbf{\gu{ઝાંખું}} (\emph{jhā̃khũ}) No dictionary we use records an older form, so its history is not settled.
\end{cousinweb}

\begin{grammarlens}[title={What you've built}]
One more word for getting around.
\end{grammarlens}

\subsection*{Your turn}
\begin{itemize}
\raggedright
  \item \emph{Say it:} \emph{jhā̃khũ}
  \item \emph{Say it:} \emph{jhā̃khũ}, once more
  \item \emph{Say it:} \emph{orḍāmā̃ jhā̃khī battī hatī.} — There was a dim light in the room.
  \item \emph{Recall:} say the Gujarati for the light of daybreak, then the Gujarati for dusk, then say \emph{jhā̃khũ} again
\end{itemize}

\begin{tcolorbox}[breakable,colback=teal!4,colframe=teal!35!black,title={Before you move on}]
What does \gu{ઝાંખું} mean? (\textquotedblleft{}Dim; faint (of light)\textquotedblright{}.) Say it once more.
\end{tcolorbox}
\section[\texorpdfstring{\gu{ધૂંધળું} (dhū̃dhḷũ)}{ધૂંધળું (dhū̃dhḷũ)}]{\gu{ધૂંધળું} (dhū̃dhḷũ) — hazy; misty}

\label{lesson:GU-C438-dhudhlu}



\begin{quote}
Before the new one: say the Gujarati for dusk, then the Gujarati for dim.
\end{quote}

\subsection*{You'll want to know: \gu{ધૂંધળું}}
\textbf{\gu{ધૂંધળું}} — \emph{dhū̃dhḷũ} — \textquotedblleft{}hazy; misty\textquotedblright{}.

\begin{quote}
\gu{ધુમ્મસમાં} \gu{બધું} \gu{ધૂંધળું} \gu{લાગે} \gu{છે}.
\end{quote}

\begin{quote}
\emph{dhummasmā̃ badhũ dhū̃dhḷũ lāge che.}
\end{quote}

\begin{quote}
In the fog everything looks hazy.
\end{quote}

\begin{cousinweb}
\textbf{\gu{ધૂંધળું}} (\emph{dhū̃dhḷũ}) goes back to \emph{*dhūmāndha}, \textquotedblleft{}blinded by smoke\textquotedblright{}, a form CDIAL reconstructs; CDIAL lists the Gujarati word, \textquotedblleft{}misty\textquotedblright{} (CDIAL 6858).
\end{cousinweb}

\begin{grammarlens}[title={What you've built}]
One more word for getting around.
\end{grammarlens}

\subsection*{Your turn}
\begin{itemize}
\raggedright
  \item \emph{Say it:} \emph{dhū̃dhḷũ}
  \item \emph{Say it:} \emph{dhū̃dhḷũ}, once more
  \item \emph{Say it:} \emph{dhummasmā̃ badhũ dhū̃dhḷũ lāge che.} — In the fog everything looks hazy.
  \item \emph{Recall:} say the Gujarati for dusk, then the Gujarati for dim, then say \emph{dhū̃dhḷũ} again
\end{itemize}

\begin{tcolorbox}[breakable,colback=teal!4,colframe=teal!35!black,title={Before you move on}]
What does \gu{ધૂંધળું} mean? (\textquotedblleft{}Hazy; misty\textquotedblright{}.) Say it once more.
\end{tcolorbox}
\section[\texorpdfstring{\gu{ટમટમવું} (ṭamṭamvũ)}{ટમટમવું (ṭamṭamvũ)}]{\gu{ટમટમવું} (ṭamṭamvũ) — to twinkle}

\label{lesson:GU-C438-tamtamvu}



\begin{quote}
Before the new one: say the Gujarati for dim, then the Gujarati for hazy.
\end{quote}

\subsection*{You'll want to know: \gu{ટમટમવું}}
\textbf{\gu{ટમટમવું}} — \emph{ṭamṭamvũ} — \textquotedblleft{}to twinkle\textquotedblright{}.

\begin{quote}
\gu{રાત્રે} \gu{આકાશમાં} \gu{તારા} \gu{ટમટમે} \gu{છે}.
\end{quote}

\begin{quote}
\emph{rātre ākāśmā̃ tārā ṭamaṭme che.}
\end{quote}

\begin{quote}
At night the stars twinkle in the sky.
\end{quote}

Past (it …): \gu{તે} \gu{તારો} \gu{ટમટમ્યો}, \emph{te tāro ṭamṭamyo}.

\begin{cousinweb}
\textbf{\gu{ટમટમવું}} (\emph{ṭamṭamvũ}) doubles a sound, as words that name a flicker often do. No dictionary we use records an older form, so its history is not settled.
\end{cousinweb}

\begin{grammarlens}[title={What you've built}]
One more word for getting around.
\end{grammarlens}

\subsection*{Your turn}
\begin{itemize}
\raggedright
  \item \emph{Say it:} \emph{ṭamṭamvũ}
  \item \emph{Say it:} \emph{ṭamṭamvũ}, once more
  \item \emph{Say it:} \emph{rātre ākāśmā̃ tārā ṭamaṭme che.} — At night the stars twinkle in the sky.
  \item \emph{Recall:} say the Gujarati for dim, then the Gujarati for hazy, then say \emph{ṭamṭamvũ} again
\end{itemize}

\begin{tcolorbox}[breakable,colback=teal!4,colframe=teal!35!black,title={Before you move on}]
What does \gu{ટમટમવું} mean? (\textquotedblleft{}To twinkle\textquotedblright{}.) Say it once more.
\end{tcolorbox}
